% !TeX program = xelatex
% UTF-8. Standalone: xelatex leukemia.tex, from this directory.
% Inclusion: define \LeukemiaInNotebook before \input in a XePersian document.
% Source: Google Drive Notebook/leukemia, 13 photographs in filename order.
% This is a transcription, not a medical update. Source wording is retained.
\ifdefined\LeukemiaInNotebook
  \let\LeukemiaFinish\relax
\else
  \def\LeukemiaFinish{\end{document}}
  \documentclass[12pt,a4paper]{article}
  \usepackage[margin=22mm]{geometry}
  \usepackage{amsmath,amssymb}
  \usepackage{graphicx}
  \usepackage{enumitem}
  \usepackage{xepersian}
  \IfFontExistsTF{Amiri}{\settextfont{Amiri}}{\settextfont{DejaVu Sans}}
  \setlatintextfont{DejaVu Serif}
  \setlength{\parindent}{0pt}
  \setlength{\parskip}{5pt}
  \setlength{\emergencystretch}{3em}
  \begin{document}
\fi
\providecommand{\LeukemiaUnclear}[1]{\textbf{[خوانش نامطمئن: #1]}}
\providecommand{\LeukemiaPage}[1]{\clearpage\section*{تصویر #1}}
\providecommand{\LeukemiaMath}[1]{\lr{\(#1\)}}
\providecommand{\LeukemiaImagePath}{images/}

\begin{center}
{\Large\bfseries لوسمی ـ ۵۴۴}\par
\lr{Nelson-22}
\end{center}
\textbf{یادداشت رونویسی:} متن فارسی و انگلیسی، حاشیه‌ها و اصلاحات خوانا به ترتیب تصاویر بازنویسی شده‌اند. موارد نامطمئن مشخص شده‌اند. تصاویر کامل در پیوست حفظ شده‌اند. سربرگ چاپی همهٔ برگه‌ها \lr{Subject:} و \lr{Date: / /} و نشان چاپی پایین صفحه \lr{YASHA} است. موضوع تصویر ۱: \lr{Nelson-22}؛ موضوع سایر تصاویر خالی است. سربرگ دست‌نویس تصاویر ۲ تا ۱۳ «۵۴۴ ـ» است؛ شماره‌های دست‌نویس پایین برگه‌ها ۱ تا ۱۳ هستند.

\LeukemiaPage{۱}
\section*{لوسمی}
لوسمی‌ها شایع‌ترین سرطان‌های کودکی هستند.
\begin{itemize}
\item ۳۰\% تمام سرطان‌ها زیر ۱۵ سال.
\item ۷۷\% تمام لوسمی‌های کودکان \lr{ALL} است؛ \lr{AML}: ۱۱\%؛ \lr{CML}: \lr{2–3\%}؛ \lr{JMML}: \lr{1–2\%} (جوونایل میلومونوسیتیک؛ یک واژهٔ میانی خط خورده است).
\end{itemize}
در لوسمی‌ها تغییرات ژنتیکی در سلول‌های هماتوپوئتیک منجر به پرولیفراسیون آنرگولیتد کلون‌های سلولی می‌شود.
\begin{itemize}
\item سلول‌ها جهش‌یافته، سرعت تکثیر بالا و نرخ آپوپتوز خودبه‌خودی پایین‌تری دارند.
\item نتیجه: قطع و اختلال عملکرد طبیعی مغز استخوان است.
\end{itemize}
\section*{\lr{ALL} (لوسمی لنفوبلاستیک حاد)}
\begin{itemize}
\item اولین سرطان دیسمنتد قابل درمان است.
\item شایع در سن \lr{2–3} سالگی.
\item در مبتلایان به برخی اختلالات ژنتیکی شایع‌تر است:
\begin{itemize}
\item \lr{Down Syndrome}
\item \lr{Bloom Syndrome}
\item \lr{Li Fraumani Syndrome} [املای دست‌نویس]
\item \lr{Ataxia-Telangiectasia}
\end{itemize}
\item برخورد با اشعهٔ \lr{X} پیش از تولد و بعد از تولد ریسک آن را بالا می‌برد.
\item در برخی کشورهای در حال توسعه بین \lr{EBV} و \lr{B-cell ALL} ارتباطاتی هست.
\end{itemize}
\subsection*{دسته‌بندی \lr{ALL}}
\begin{itemize}
\item وابسته به ویژگی‌های سلول‌های بدخیم در مغز استخوان است.
\item مورفولوژی سلول‌ها برای تشخیص \lr{ALL} کافی است.
\end{itemize}

\LeukemiaPage{۲}
۸۵\% تمام \lr{ALL}ها از نوع \lr{B-ALL} و ۱۵\% از نوع \lr{T-ALL} هستند و البته نزدیک به ۱\% نیز از نوع تمایزگرفتهٔ \lr{mature B cells} هستند.
\begin{itemize}
\item به نوع آخر \lr{Burkitt Leukemia} می‌گویند.
\item خیلی سریع رشد می‌کند و نیازمند درمان متفاوت است.
\end{itemize}
در دسته‌ای از کودکان نوعی از لوسمی را داریم که سلول‌های آن مارکرهای هر دو نوع لنفوئید و میلوئید را بیان می‌کنند.

از اختلالات کروموزومی برای \lr{sub-classify} کردن لوسمی‌ها از نظر پروگنوز استفاده می‌شود.

در دورهٔ فالوآپ پس از درمان، ما دنبال \lr{MRD: Minimal Residual Disease}. [جمله در دست‌نویس در همین‌جا پایان می‌یابد.]
\section*{تظاهرات بالینی}
\begin{itemize}
\item تظاهرات اولیه معمولاً غیراختصاصی و کوتاه هستند.
\item آنورکسی، خستگی، ضعف، تحریک‌پذیری، تب \lr{low grade} اینترمیتنت.
\item درد عضلانی و استخوانی خصوصاً در اندام تحتانی.
\item تورم مفاصل.
\item درد استخوانی معمولاً شدید است و کودک را از خواب بیدار می‌کند.
\item به دنبال پیشرفت بیماری شواهد نارسایی \lr{BM} نمایان می‌شود:
\begin{itemize}
\item رنگ‌پریدگی، خستگی، عدم تحمل فعالیت، کبودی، خونریزی مخاطی، تب (به دلیل لوسمی و عفونت).
\end{itemize}
\item لنفادنوپاتی، هپاتواسپلنومگالی، بزرگ شدن بیضه.
\item درگیری \lr{CNS}: نوروپاتی کرانیال، سردرد، تشنج.
\end{itemize}

\LeukemiaPage{۳}
\section*{تظاهرات بالینی (ادامه)}
\begin{itemize}
\item دیسترس تنفسی، آپنه، فشار تودهٔ مدیاستینال بر راه‌های هوایی.
\item معاینهٔ فیزیکی:
\begin{itemize}
\item رنگ‌پریدگی، پتشی، پورپورا، خونریزی مخاطی.
\item لنفادنوپاتی شایع‌تر از هپاتومگالی است.
\item تندرنس استخوان و مفصل.
\item کودک درد استخوان عمقی داشته باشد، بدون تندرنس است.
\end{itemize}
\item شواهد \LeukemiaMath{\uparrow ICP}: نشان‌دهندهٔ درگیری لوکمیک \lr{CNS} است.
\begin{itemize}
\item پاپیل‌ادما؛ پارزی عصب کرانیال.
\item خونریزی رتینال.
\end{itemize}
\item وجود یک تودهٔ بزرگ مدیاستن (تیموسی یا نود).
\begin{itemize}
\item معمولاً در نوجوان پسر با \lr{T-ALL} شایع‌تر است.
\end{itemize}
\item \lr{T-ALL} معمولاً با تعداد بیشتر لکوسیت همراه است.
\end{itemize}
\section*{\lr{B-ALL}}
\begin{itemize}
\item شایع‌ترین ایمونوفنوتیپ \lr{ALL} است.
\item شروع بیماری در \lr{1–10} سالگی.
\item تعداد میانهٔ لکوسیت ۳۳٬۰۰۰ است؛ ۷۵\% زیر ۲۰٬۰۰۰.
\item در ۷۵\% موارد ترومبوسایتوپنی دیده می‌شود.
\item در \lr{30–40\%} موارد هپاتواسپلنومگالی دیده می‌شود.
\item در ۵\% درگیری \lr{CNS} دیده می‌شود. [\lr{CSF} خط خورده و \lr{CNS} جایگزین شده است.]
\item در \lr{10–15\%} شواهدی در \lr{CSF} دیده می‌شود ولی تنها در ۵\% برای \lr{CNS Leukemia} تشخیصی است.
\begin{itemize}
\item \LeukemiaMath{WBC>5} در هر میکرولیتر، همراه با بلاست.
\end{itemize}
\item ۲۵\% پسران درگیری بیضه.
\end{itemize}

\LeukemiaPage{۴}
اندیکاسیون بیوپسی \lr{testis} وجود ندارد.
\section*{تشخیص لوسمی \lr{ALL}}
\begin{itemize}
\item بر اساس \lr{PBS} است؛ که نشان‌دهندهٔ نارسایی مغز استخوان است.
\item آنمی و ترومبوسایتوپنی در اغلب بیماران.
\item بسیاری از بیماران \lr{ALL} لکوسیت کمتر از ۱۰٬۰۰۰ دارند.
\item در سنین پایین‌تر، سلول‌های لوکمیک اغلب در ابتدا لنفوسیت آتیپیک گزارش می‌شوند.
\item زمانی که \lr{PBS} مطرح‌کنندهٔ احتمال بالای لوسمی است، فلوسایتومتری محیطی می‌تواند به سرعت تشخیص قطعی را مطرح کند.
\item برای تأیید تشخیص \lr{BMB} نیاز است:
\begin{itemize}
\item \LeukemiaMath{>25\%} سلول‌های مغز استخوان جمعیت هموژنوسی از لنفوبلاست باشند.
\end{itemize}
\item بررسی اولیه شامل \lr{CSF} نیز می‌شود.
\item اگر سلول بلاست دیده شد و لکوسیت \lr{CSF} نیز \LeukemiaMath{\uparrow} بود، نشانهٔ لوکمیا مننژ است.
\begin{itemize}
\item مطرح‌کنندهٔ پروگنوز ضعیف‌تر و نیاز به درمان مضاعف \lr{CNS}.
\end{itemize}
\item \lr{Staging LP} ممکن است در همراهی با اولین دوز کموتراپی اینتراتکال انجام شود.
\item فرد توانمند باید \lr{LP} را انجام دهد چون که یک \lr{LP} تروماتیک ممکن است خطر عود \lr{CNS} را بالا ببرد.
\end{itemize}

\LeukemiaPage{۵}
\section*{تشخیص‌های افتراقی \lr{ALL}}
\begin{itemize}
\item علائم تیپیک + نشانه‌ها + آنمی + ترومبوسایتوپنی + لکوسیتوز + بلاست در \lr{PBS}: تشخیص تقریباً قطعی است.
\item \lr{LDH} معمولاً در \lr{ALL} بالا است.
\end{itemize}
\begin{enumerate}
\item اگر پان‌سایتوپنی بود، بدون بلاست در \lr{PBS}:
\begin{itemize}
\item آنمی آپلاستیک (کانستیتوشنال یا اکتسابی)، درگیری استخوان ناشی از بیماری متاستاتیک، \lr{HLH} (هموفاگوسیتیک لنفوهیستیوسیتوزیس).
\end{itemize}
\item اگر یک لاین سلولی افت کند:
\begin{itemize}
\item \LeukemiaUnclear{«ترومبوبلاستوپنی گذرا کودکان»؛ واژهٔ اول ممکن است «اریتروبلاستوپنی» باشد}، ترومبوسایتوپنی ایمنی و نوتروپنی مادرزادی یا اکتسابی.
\end{itemize}
\end{enumerate}
۱ و ۲ نیز می‌توانند تظاهر اولیهٔ \lr{ALL} باشند.
\begin{itemize}
\item تب حاد + لنفادنوپاتی: \lr{ALL}؛ \lr{IMN}.
\item تب + درد استخوانی: \lr{ALL}؛ \lr{JIA}.
\item سرطان‌های زیر با همین آزمایش‌ها ممکن است شبیه \lr{ALL} باشند:
\begin{itemize}
\item \lr{AML}، نوروبلاستوما، رابدومیوسارکوما، \lr{Ewing Sarcoma}، رتینوبلاستوما.
\end{itemize}
\end{itemize}
\section*{درمان \lr{ALL}}
\begin{itemize}
\item مهم‌ترین فاکتور پروگنوستیک، درمان \lr{ALL} است.
\item بدون درمان مناسب \lr{ALL} کشنده است.
\item رویکرد استاندارد مدرن \lr{ALL}: \lr{Risk Stratified Therapy} است.
\item ریسک‌های \lr{high risk}:
\begin{itemize}
\item سن زیر ۱ سال، بالای ۱۰ سال.
\item لکوسیت اولیه بالای ۵۰٬۰۰۰. [متن روی نشان چاپی \lr{YASHA} نوشته شده است.]
\item جنسیت پسر. [یادداشت حاشیهٔ پایین]
\end{itemize}
\end{itemize}

\LeukemiaPage{۶}
\section*{درمان \lr{ALL} (ادامه)}
\begin{itemize}
\item درمان استاندارد شامل \lr{ChT} برای \lr{2–3} سال است.
\item درمان اولیه \lr{Remission Induction} نام دارد.
\item هدف: پاک‌سازی مغز استخوان از سلول لوکمیک است.
\begin{itemize}
\item \lr{Vincristine} هفتگی برای ۴ هفته.
\item یک کورتون مثل دگزا یا پردنیزولون.
\item یک دوز طولانی‌اثری از یک \lr{Pegylated Asparaginase}.
\item \lr{Daunorubicin} هفتگی در موارد \lr{high risk}.
\end{itemize}
\item با این اپروچ ۹۸\% موارد وارد رمیشن می‌شوند که به معنای بلاست کمتر از ۵\% است و بازگشت تعداد نوتروفیل و \lr{Plt} به نزدیک نرمال بعد از \lr{4–5} هفتهٔ درمان.
\item \lr{ChT} اینتراتکال یک دوز در ابتدا و حداقل یک دوز دیگر در طول دوره داده می‌شود.
\end{itemize}
\subsection*{فاز دوم درمان: \lr{Consolidation}}
\begin{itemize}
\item درمان اینتنسیو \lr{CNS} همراه با درمان سیستمیک جهت جلوگیری از عود \lr{CNS}.
\item \lr{ChT} اینتراتکال مکرر توسط \lr{LP} انجام می‌شود.
\item ریسک عود \lr{CNS} به \LeukemiaMath{<5\%} کاهش می‌یابد.
\item اغلب رژیم‌ها \lr{14–28} هفته درمان دارند.
\end{itemize}
\subsection*{فاز سوم: \lr{Maintenance}}
\lr{2–3} سال.
\begin{itemize}
\item مرکاپتوپورین روزانه.
\item متوترکسات هفتگی خوراکی.
\item دوزهای اینترمیتنت \lr{Vincristine} و کورتیکواستروئید.
\end{itemize}

\LeukemiaPage{۷}
نوجوانان و جوانان کم‌سن با \lr{ALL} پروگنوز ضعیف‌تری از افراد \LeukemiaMath{<15} سال دارند.
\begin{itemize}
\item عود در مغز استخوان \lr{15–20\%} از موارد \lr{ALL} رخ می‌دهد.
\item ریسک عود در \lr{testis} کمتر از ۲\% است.
\begin{itemize}
\item عود به صورت تورم بدون درد یک یا دو بیضه.
\end{itemize}
\end{itemize}
بیماران با \lr{WBC} بالا مستعد سندرم لیز تومور هستند.

نارسایی کلیه در همراهی با اوریک‌اسید خیلی بالا ممکن است دیده شود که برای پیشگیری از آن، از آلوپورینول و اورات اکسیداز استفاده می‌کنیم.

شیمی‌درمانی معمولاً سرکوب مغز استخوان ایجاد می‌کند و بیمار معمولاً نیازمند تزریق \lr{P.C} و \lr{Plt} می‌شود.

در کودک تب‌دار با نوتروپنی باید خیلی به سپسیس حساس بود و آنتی‌بیوتیک وسیع‌الطیف تجربی را شروع کرد.

بیماران باید حین درمان و تا ماه‌ها بعد از آن به طور پروفیلاکسی برای پنومونی ناشی از پنوموسیستیس جیرووسی \lr{ABT} دریافت کنند.

درمان بیماری باید شدید باشد. این درمان شدید ممکن است ارگان‌های بدن مثل \LeukemiaUnclear{«غدد»} را نیز به طور دائمی آسیب بزنند.

یادداشت حاشیهٔ پایین: \LeukemiaUnclear{«بلاکرهای T ... پروتکشن ...»؛ پایان عبارت زیر نشان چاپی YASHA قرار گرفته است}.

\LeukemiaPage{۸}
\section*{پیش‌آگهی \lr{ALL}}
\begin{itemize}
\item \lr{5YS} کلی آن ۹۰\% است.
\item بیمارانی که \lr{Cure} می‌شوند، با احتمال بالا آسیب‌های قلبی ناشی از درمان را تجربه خواهند کرد.
\end{itemize}
\section*{\lr{AML}}
\subsection*{اپیدمیولوژی}
\begin{itemize}
\item ۱۱\% لوسمی‌های کودکان / ۳۷\% لوسمی‌های \lr{15–19} سال.
\item \lr{APL} (لوسمی پرومیلوسیتیک حاد) یک ساب‌تایپ \lr{AML} است که در برخی مناطق جغرافیایی شایع‌تر است.
\end{itemize}
\subsection*{ریسک فاکتورها}
\begin{itemize}
\item اشعهٔ یونیزه.
\item داروهای کموتراپی (\lr{Alkylating agent}، اپی‌پودوفیلوتوکسین).
\item حلال‌های ارگانیک (\lr{benzene}).
\item جهش‌های قبلی (در ۱۰\% موارد \lr{AML}).
\item برخی سندرم‌ها:
\begin{itemize}
\item سندرم داون.
\item آنمی فانکونی.
\item سندرم \lr{Bloom}.
\item سندرم \lr{Kostmann}.
\item سندرم \lr{Shwachman-Diamond}.
\item سندرم \lr{Diamond-Blackfan}.
\item سندرم \lr{Li-Fraumeni}.
\item نوروفیبروماتوز تیپ \lr{I}.
\end{itemize}
\end{itemize}

\LeukemiaPage{۹}
\section*{طبقه‌بندی سلولی \lr{AML}}
\begin{itemize}
\item ویژگی کاراکتریستیک: \LeukemiaMath{>20\%} مغز استخوان، جمعیت هموژنوسی از سلول‌های بلاست باشد با ویژگی‌های مشابه آن‌هایی که تمایز به ردهٔ میلوئید، مونوسیت، مگاکاریوسیت را نمایش می‌دهند.
\end{itemize}
\section*{تظاهرات بالینی}
\begin{itemize}
\item یافتهٔ مشترک: جایگزینی مغز استخوان با سلول‌های بدخیم است.
\item بیماران \lr{AML} می‌توانند با ویژگی‌های \lr{ALL} پرزنت کنند.
\item یا با ویژگی‌هایی که در \lr{ALL} نادر است:
\begin{itemize}
\item ندول‌ها (ساب‌کوتانئوس) یا ضایعات بلوبری مافین (خصوصاً در شیرخواران).
\item هیپرتروفی لثه‌ها (خصوصاً در ساب‌تایپ مونوسیت).
\item شواهد آزمایشگاهی \lr{DIC} (جهت مطرح کردن \lr{APL} می‌باشد).
\item توده‌های مجزا (کلوروما یا گرانولوسیتیک سارکوما):
\begin{itemize}
\item می‌تواند در فقدان سایر شواهد مغز استخوان رخ دهد.
\item با ترانسلوکیشن \lr{t(8,21)} مرتبط است.
\end{itemize}
\end{itemize}
\item کلوروما: در اوربیت یا فضای اپیدورال دیده می‌شود.
\end{itemize}
\section*{تشخیص \lr{AML}}
بیوپسی مغز استخوان: مغز استخوان هایپرسلولار با الگوی سلولی مونوتون.
\section*{پروگنوز و درمان \lr{AML}}
\begin{itemize}
\item شیمی‌درمانی چنددارویی تهاجمی در \lr{85–90\%} موارد بیمار را به رمیشن \lr{C.R} \LeukemiaUnclear{حاشیهٔ چپ پس از «رمیشن»، شبیه C.R}.
\item نرخ سروایول \lr{60–70\%} است.
\end{itemize}

\LeukemiaPage{۱۰}
\section*{درمان \lr{AML} (ادامه)}
\begin{itemize}
\item رژیم اینداکشن: \lr{Anthracycline} + سیتارابین \lr{high dose}.
\item ۵\% بیماران قبل از رسیدن به اینداکشن در اثر خونریزی یا عفونت فوت می‌کنند.
\item \lr{APL}: بیمار به درمان زیر خوب جواب می‌دهد:
\begin{latin}
\noindent all-trans-retinoic acid (ATRA, tretinoin) + anthracycline + cyta
\end{latin}
\item نرخ مرگ بر اثر عفونت بسیار بالا است.
\begin{itemize}
\item استرپ وریدانس؛ قارچ.
\end{itemize}
\end{itemize}
\section*{سندرم \lr{Down}، لوسمی حاد و \lr{Transient Abnormal Myelopoiesis}}
\begin{itemize}
\item لوسمی حاد با احتمال \lr{15–20} برابر در سندرم داون رخ می‌دهد.
\item نسبت \lr{ALL} به \lr{AML} در سندرم داون مشابه مردم عادی است.
\begin{itemize}
\item بجز ۳ سال اول زندگی که \lr{AML} شایع‌تر است.
\end{itemize}
\item پروگنوز لوسمی در سندرم داون بهتر است.
\end{itemize}
۱۰\% نوزادان با سندرم داون دچار \lr{TAM (Transient Abnormal Myelopoiesis)} می‌شوند.
\begin{itemize}
\item \LeukemiaMath{\uparrow WBC}، بلاست در \lr{PBS}، آنمی، ترومبوسایتوپنی، زردی، هپاتواسپلنومگالی.
\item هیدروپس جنینی یک تظاهر نامعمول \lr{TAM} است؛ در همراهی با پلورال افیوژن، آسیت، نارسایی قلبی، اختلال عملکرد کبدی، آنمی.
\item مرگ بر اثر \lr{TAM} نادر است ولی در زمینهٔ هایپرلکوسیتوز، آسیت، کوآگولوپاتی، پره‌مچوریتی، نارسایی کبدی و کلیوی رخ می‌دهد.
\item علائم معمولاً ظرف ۳ ماه اول زندگی رفع می‌شود.
\end{itemize}

\LeukemiaPage{۱۱}
\section*{درمان \lr{TAM}}
\begin{itemize}
\item نیاز به حمایت دارند ولی \lr{ChT} نه.
\item سیتارابین با دوز پایین در افراد با \LeukemiaMath{\uparrow WBC}.
\end{itemize}
\lr{20–30\%} افرادی که دچار \lr{TAM} شده‌اند دچار لوسمی کلینیکال می‌شوند؛ ظرف ۳ سال، متوسط ۱۶ ماه.
\begin{itemize}
\item معمولاً \lr{Acute Megakaryocytic leukemia}.
\end{itemize}
\section*{\lr{CML: Chronic Myelogenous Leukemia}}
\lr{2–3\%} تمام موارد لوسمی در کودکان را شامل می‌شود.

۹۹\% موارد \lr{CML} ترانسلوکیشن \lr{t(9,22)(q34,q11)} معروف به کروموزوم فیلادلفیا را دارند؛ پروتئین \lr{BCR-ABL fusion}.
\begin{itemize}
\item در مواردی از \lr{ALL} هم دیده می‌شود.
\end{itemize}
\subsection*{علائم بالینی \lr{CML}}
غیراختصاصی هستند.
\begin{itemize}
\item تب، خستگی، کاهش وزن، آنورکسی.
\item اسپلنومگالی؛ درد \lr{RUQ}. [همین اختصار در دست‌نویس آمده است.]
\end{itemize}
\subsection*{تست تشخیصی}
لکوسیتوز با سلول‌های میلوئید در تمام استیج‌های تمایز در \lr{PBS}.
\subsection*{فاز \lr{Chronic} بیماری}
\begin{itemize}
\item \LeukemiaMath{\uparrow WBC}، آنمی خفیف، ترومبوسایتوپنی.
\item \lr{3–4} سال طول می‌کشد.
\end{itemize}
\subsection*{فاز \lr{Accelerated (blast crisis)}}
\begin{itemize}
\item تعداد سلول‌ها شدیداً بالا می‌رود.
\item تظاهر بالینی از لوسمی حاد غیرقابل افتراق است.
\end{itemize}

\LeukemiaPage{۱۲}
\section*{درمان \lr{CML}}
ایماتینیب: مهارکنندهٔ اختصاصی \lr{BCR-ABL} تیروزین کیناز.

نسل دوم تیروزین کیناز اینهیبیتورها: نرخ رمیشن بالاتر.
\begin{itemize}
\item دازاتینیب، نیلوتینیب.
\end{itemize}
هیدروکسی‌اوره، \lr{CML} را در فاز مزمن کنترل می‌کند و تعداد \lr{WBC} را نرمال می‌کند.

درمان استاندارد \lr{CML} کودکان در حال حاضر مهارکنندهٔ تیروزین کیناز است.
\section*{\lr{JMML: Juvenile Myelomonocytic Leukemia}}
\begin{itemize}
\item قبلاً به آن \lr{CML} جوانان می‌گفتند.
\item معمولاً مربوط به \LeukemiaMath{<2} سال است.
\item ۱\% تمام لوسمی‌های کودکان را شامل می‌شود.
\item کروموزوم فیلادلفیا ندارند.
\end{itemize}
\subsection*{علائم بالینی}
\begin{itemize}
\item راش، لنفادنوپاتی، اسپلنومگالی، خونریزی.
\item \LeukemiaMath{\uparrow WBC}، \LeukemiaMath{\uparrow} مونوسیت، ترومبوسایتوپنی، آنمی، اریتروبلاست.
\end{itemize}
در مغز استخوان: یک الگوی میلودیسپلاستیک؛ بلاست \LeukemiaMath{<20\%}.

\lr{RAS Oncogenic path} فعال می‌شود در \lr{JMML}.

\LeukemiaPage{۱۳}
\section*{\lr{JMML} (ادامه)}
بیماران \lr{NF-1} و سندرم \lr{Noonan} مستعد \lr{JMML} هستند.
\begin{itemize}
\item درمان \lr{JMML}: به طور کلی پیوند مغز استخوان؛ پروگنوز همچنان \lr{poor}.
\item در همراهی سندرم \lr{Noonan} ممکن است خودبه‌خود بهبود یابد.
\end{itemize}
\section*{\lr{Infant Leukemia}}
\begin{itemize}
\item \LeukemiaUnclear{۲ یا ۳}\% تمام لوسمی‌ها قبل از ۱ سالگی.
\item هایپرلکوسیتوز + درگیری بافتی شدید: ارگانومگالی، درگیری \lr{CNS}.
\item ندول ساب‌کوتانئوس (لوکمیا کوتیس) + فلج عصب (به دنبال درگیری \lr{CNS}).
\end{itemize}

\clearpage
\section*{پیوست: تصاویر کامل منبع}
هر تصویر با نام اصلی و شناسهٔ فایل گوگل درایو در توضیحات کد ثبت شده است. برای استفاده از این فایل با \lr{\texttt{\string\input}} از مسیر دیگری، دستور \lr{\texttt{\string\LeukemiaImagePath}} را به مسیر پوشهٔ تصاویر تنظیم کنید.
% pic01.jpg = 20261010_051318530.jpg; Drive ID: 1E5YEXYdNQ5-tnSOEaBFqSWzAoulc5i-P
\clearpage
\subsection*{تصویر 1: \lr{pic01.jpg}}
\begin{center}
\includegraphics[width=\linewidth,height=0.87\textheight,keepaspectratio]{\LeukemiaImagePath pic01.jpg}
\end{center}
% pic02.jpg = 20261010_051322484.jpg; Drive ID: 1dHLAWHpk_cSRnod3NYmY9rNr79rvpDRp
\clearpage
\subsection*{تصویر 2: \lr{pic02.jpg}}
\begin{center}
\includegraphics[width=\linewidth,height=0.87\textheight,keepaspectratio]{\LeukemiaImagePath pic02.jpg}
\end{center}
% pic03.jpg = 20261010_051329481.jpg; Drive ID: 1DQr4XOwCw2o_4_phuXbn_8kAuHdmcW36
\clearpage
\subsection*{تصویر 3: \lr{pic03.jpg}}
\begin{center}
\includegraphics[width=\linewidth,height=0.87\textheight,keepaspectratio]{\LeukemiaImagePath pic03.jpg}
\end{center}
% pic04.jpg = 20261010_051334015.jpg; Drive ID: 15lacjsriTIKt86MhIGeKaiNYNn7gxuhx
\clearpage
\subsection*{تصویر 4: \lr{pic04.jpg}}
\begin{center}
\includegraphics[width=\linewidth,height=0.87\textheight,keepaspectratio]{\LeukemiaImagePath pic04.jpg}
\end{center}
% pic05.jpg = 20261010_051343326.jpg; Drive ID: 1VQ1M8-E6vhBV0rR46SHqV5YosquizFbn
\clearpage
\subsection*{تصویر 5: \lr{pic05.jpg}}
\begin{center}
\includegraphics[width=\linewidth,height=0.87\textheight,keepaspectratio]{\LeukemiaImagePath pic05.jpg}
\end{center}
% pic06.jpg = 20261010_051348604.jpg; Drive ID: 1Qih-a0_6atU0vxooD6Mh97fTL608lg7R
\clearpage
\subsection*{تصویر 6: \lr{pic06.jpg}}
\begin{center}
\includegraphics[width=\linewidth,height=0.87\textheight,keepaspectratio]{\LeukemiaImagePath pic06.jpg}
\end{center}
% pic07.jpg = 20261010_051357033.jpg; Drive ID: 1DnSLI9x9hu1rZFdTemcL1YIUaN8dF-5Q
\clearpage
\subsection*{تصویر 7: \lr{pic07.jpg}}
\begin{center}
\includegraphics[width=\linewidth,height=0.87\textheight,keepaspectratio]{\LeukemiaImagePath pic07.jpg}
\end{center}
% pic08.jpg = 20261010_051401420.jpg; Drive ID: 1cM-fAqJQ7ksRzXriJnle7FnurkEDfLEw
\clearpage
\subsection*{تصویر 8: \lr{pic08.jpg}}
\begin{center}
\includegraphics[width=\linewidth,height=0.87\textheight,keepaspectratio]{\LeukemiaImagePath pic08.jpg}
\end{center}
% pic09.jpg = 20261010_051411878.jpg; Drive ID: 1ZJg8WP7o_07JsHR016uTWCucHg4G1WYT
\clearpage
\subsection*{تصویر 9: \lr{pic09.jpg}}
\begin{center}
\includegraphics[width=\linewidth,height=0.87\textheight,keepaspectratio]{\LeukemiaImagePath pic09.jpg}
\end{center}
% pic10.jpg = 20261010_051415617.jpg; Drive ID: 1hNbKs3j-JycY2bFomS_HTm54egzTqNWx
\clearpage
\subsection*{تصویر 10: \lr{pic10.jpg}}
\begin{center}
\includegraphics[width=\linewidth,height=0.87\textheight,keepaspectratio]{\LeukemiaImagePath pic10.jpg}
\end{center}
% pic11.jpg = 20261010_051423010.jpg; Drive ID: 1bGpMJGBaycVXJjeFJM9RqdtMxJy6tvyN
\clearpage
\subsection*{تصویر 11: \lr{pic11.jpg}}
\begin{center}
\includegraphics[width=\linewidth,height=0.87\textheight,keepaspectratio]{\LeukemiaImagePath pic11.jpg}
\end{center}
% pic12.jpg = 20261010_051426747.jpg; Drive ID: 1ZCYIW2U0mdIjP9P-6rcY70Dw4X7zcCsL
\clearpage
\subsection*{تصویر 12: \lr{pic12.jpg}}
\begin{center}
\includegraphics[width=\linewidth,height=0.87\textheight,keepaspectratio]{\LeukemiaImagePath pic12.jpg}
\end{center}
% pic13.jpg = 20261010_051433443.jpg; Drive ID: 1sSzJS9clPNQj-fkn4HGw6lruCjvmDxF0
\clearpage
\subsection*{تصویر 13: \lr{pic13.jpg}}
\begin{center}
\includegraphics[width=\linewidth,height=0.87\textheight,keepaspectratio]{\LeukemiaImagePath pic13.jpg}
\end{center}
\LeukemiaFinish
